\section*{अध्याय \ref{ch:b1:p-block-16-18} --- p-ब्लॉक II: ऑक्सीजन, हैलोजन और उत्कृष्ट गैसें}

\begin{solution}{exo:b1:p-block-16-18:1}
\ce{H2S} $-$II; \ce{S8} 0; \ce{SO2} $+$IV; \ce{SO3^2-} $+$IV; \ce{SO4^2-}
$+$VI; \ce{S2O3^2-} औसतन $+$II; \ce{H2S2O7} $+$VI।
\end{solution}

\begin{solution}{exo:b1:p-block-16-18:2}
\omterm{def:b1:p-block-16-18:halogen}{हैलोजन} अपने नीचे के हैलाइडों को ऑक्सीकृत करता है: \ce{Cl2 + 2I- -> 2Cl- + I2} और
\ce{Cl2 + 2Br- -> 2Cl- + Br2} होती हैं; ब्रोमाइड के साथ आयोडीन और क्लोराइड के साथ
ब्रोमीन नहीं।
\end{solution}

\begin{solution}{exo:b1:p-block-16-18:3}
\ce{SO2}: \omterm{def:b1:lewis-resonance-vsepr:lewis}{एकाकी युग्म} वाला S, दोनों \omterm{def:b1:lewis-resonance-vsepr:resonance}{अनुनादी संरचनाओं} में से हर एक में एक S=O और एक \ce{S+-O-}
(या विस्तारित अष्टक के साथ दो S=O): कोणीय।
\ce{SO3}: त्रिकोणीय समतलीय, तीन तुल्य S--O आबंध। \ce{O3}: \omterm{def:b1:lewis-resonance-vsepr:lewis}{एकाकी युग्म} वाला केंद्रीय O,
एक O=O और एक O--O$^-$, दो \omterm{def:b1:lewis-resonance-vsepr:resonance}{अनुनादी संरचनाएँ}: कोणीय।
\end{solution}

\begin{solution}{exo:b1:p-block-16-18:4}
\ce{SO2 + H2O <=> H2SO3}; \ce{SO2 + 2NaOH -> Na2SO3 + H2O};
\ce{SO3 + H2O -> H2SO4}; \ce{SO3 + 2NaOH -> Na2SO4 + H2O}।
\end{solution}

\begin{solution}{exo:b1:p-block-16-18:5}
HF: $x^2/(0.10 - x) = 10^{-3.16}$, $x = \qty{8.0e-3}{mol/L}$, pH 2.10। HCl:
pH 1.00। H--F आबंध प्रबल है और फ़्लुओराइड आयन प्रबल \omterm{def:b1:intermolecular-forces-solvents:hbond}{हाइड्रोजन
आबंधों} में बँधा है: HF एक \omterm{def:b1:predominant-reaction:strong-acid}{दुर्बल अम्ल} है।
% ledger: dfg:HF_ao, dfg:F-_ao
\end{solution}

\begin{solution}{exo:b1:p-block-16-18:6}
\ce{BrF5}: पाँच आबंध और एक \omterm{def:b1:lewis-resonance-vsepr:lewis}{एकाकी युग्म}, वर्ग पिरामिडी। \ce{IF7}:
पंचकोणीय द्विपिरामिडी। \ce{ICl4-}: चार आबंध और दो \omterm{def:b1:lewis-resonance-vsepr:lewis}{एकाकी युग्म}, वर्ग
समतलीय। \ce{XeO3}: तीन आबंध और एक \omterm{def:b1:lewis-resonance-vsepr:lewis}{एकाकी युग्म}, त्रिकोणीय पिरामिडी।
\end{solution}

\begin{solution}{exo:b1:p-block-16-18:7}
$\log K = 2(1.36 - 0.53)/0.059 = 28$। मुक्त आयोडीन स्टार्च के साथ गहरा
नीला \omterm{def:b1:complexation:complex}{संकुल} बनाती है।
\end{solution}

\begin{solution}{exo:b1:p-block-16-18:8}
$E^\circ(\ce{O3}/\ce{O2}) = 2.07 > 0.53$~V: \ce{O3 + 2I- + 2H+ -> O2 + I2
+ H2O}। बनी आयोडीन का थायोसल्फ़ेट से \omterm{def:b1:titration-methods:titration}{अनुमापन} किया जाता है
(\cref{exo:b1:nernst:7}): एक \ce{I2}, प्रति \ce{O3}।
\end{solution}

\begin{solution}{exo:b1:p-block-16-18:9}
\ce{C12H22O11 -> 12C + 11H2O}; $M = \qty{342.0}{g/mol}$, $10.0/342.0 =
\qty{0.0292}{mol}$, कार्बन $0.0292 \times 12 \times 12.0 = \qty{4.21}{g}$।
\end{solution}

\begin{solution}{exo:b1:p-block-16-18:10}
$x(0.10 + x)/(0.10 - x) = 10^{-1.99}$: $x = \qty{8.6e-3}{mol/L}$,
$[\ce{H3O+}] = 0.1086$, 1.00 के बजाय $\mathrm{pH} = 0.96$: दूसरी
अम्लता $[\ce{H3O+}]$ में लगभग \qty{9}{\percent} जोड़ती है।
\end{solution}

\begin{solution}{exo:b1:p-block-16-18:11}
\ce{F2}/\ce{F-} (2.89~V) हर अन्य युग्म से ऊपर है: कोई भी \omterm{def:b1:nernst:couple}{ऑक्सीकारक}
फ़्लुओराइड से इलेक्ट्रॉन नहीं ले सकता। जलीय विद्युत-अपघटन में जल, जो
कहीं निम्नतर विभव (\ce{O2}/\ce{H2O}, 1.23~V) पर ऑक्सीकृत होता है, उसके बजाय अभिक्रिया करेगा;
गलित फ़्लुओराइड में कोई जल नहीं होता।
\end{solution}

\begin{solution}{exo:b1:p-block-16-18:12}
$\Delta_r G^\circ = -51.4 - (16.40 - 51.57) = \qty{-16.2}{kJ/mol}$,
$\log K = 16.2/5.708 = 2.84$, $K = 7 \times 10^2$।
% ledger: dfg:I2_ao, dfg:I-_ao, dfg:I3-_ao
\end{solution}

\begin{solution}{pb:b1:p-block-16-18:1}
\textbf{1.} 0, $+$IV, $+$VI, $+$VI।
\textbf{2.} \ce{S + O2 -> SO2}।
\textbf{3.} दो आबंधी क्षेत्रों और एक \omterm{def:b1:lewis-resonance-vsepr:lewis}{एकाकी युग्म} वाला S: कोणीय, लगभग
\ang{120}।
\textbf{4.} $10^6/32.1 = \qty{3.12e4}{mol}$।
\textbf{5.} \ce{O2} के \qty{3.115e4}{mol}: $V = 3.115 \times 10^4 \times
8.314 \times 298.15/10^5 = \qty{772}{m^3}$, वायु $772/0.21 = \qty{3.7e3}{m^3}$।
\textbf{6.} दहन अत्यंत ऊष्माक्षेपी है; \omterm{def:b1:elementary-steps:catalyst}{उत्प्रेरक} एक सीमित ताप-परास में
काम करता है, और \ce{SO2} का ऑक्सीकरण, जो स्वयं भी
ऊष्माक्षेपी है, गरम अवस्था में कम पूर्ण होता है।
\textbf{7.} \ce{2SO2 + O2 <=> 2SO3}।
\textbf{8.} बड़ा साम्य स्थिरांक दर के बारे में कुछ नहीं बताता:
कमरे के ताप पर अभिक्रिया आगे नहीं बढ़ती।
\textbf{9.} वह कम \omterm{def:b1:rate-laws:activation-energy}{सक्रियण ऊर्जा} वाला मार्ग प्रदान करता है और
पुनर्जनित होता है: एक विषमांगी \omterm{def:b1:elementary-steps:catalyst}{उत्प्रेरक}।
\textbf{10.} त्रिकोणीय समतलीय।
\textbf{11.} अधिक ऑक्सीजन भागफल
$Q = p_{\ce{SO3}}^2/(p_{\ce{SO2}}^2 p_{\ce{O2}})$ को नीचा रखती है: \ce{SO2} का
रूपांतरण और आगे जाता है।
\textbf{12.} \qty{0.3}{\percent}।
\textbf{13.} \ce{SO3 + H2SO4 -> H2S2O7}।
\textbf{14.} वह अम्ल की महीन बूँदों की ऐसी धुंध बनाता है जो
अवशोषक से निकल जाती है।
\textbf{15.} \ce{H2S2O7 + H2O -> 2H2SO4}।
\textbf{16.} \ce{2S + 3O2 + 2H2O -> 2H2SO4}।
\textbf{17.} प्रति सल्फ़र एक जल: $\num{3.115e4} \times 18.0 =
\qty{561}{kg}$।
\textbf{18.} तनुकरण बहुत ऊष्मा मुक्त करता है; अम्ल पर डाला गया जल
तुरंत उबलकर अम्ल के छींटे उड़ाता है।
\textbf{19.} $x(0.050 + x)/(0.050 - x) = 10^{-1.99}$: $x = \qty{7.5e-3}{mol/L}$,
$[\ce{H3O+}] = \qty{0.0575}{mol/L}$, $\mathrm{pH} = 1.24$।
\textbf{20.} $3.12 \times 10^4 \times 0.995 \times 98.1 = \qty{3.04}{t}$।
\textbf{21.} $84 \times 98.1/32.1 = \qty{257}{Mt}$।
\textbf{22.} $98.1/32.1 =$ \textbf{प्रति टन सल्फ़र \qty{3.06}{t} सल्फ़्यूरिक
अम्ल}।
\end{solution}
